\documentclass[english,twoside, openright, hidelinks, 11 pt, class=report,crop=false]{standalone}
\input{../bm}
\input{../../preamb}

\begin{document}

\section{Å finne størrelser}
Ligninger, formler og funksjoner er begreper som dukker opp i forskjellige sammenhenger, men som i bunn og grunn handler om det samme; \textsl{de uttrykker relasjoner mellom størrelser}. De fleste regelboksene i denne boka inneholder en formel. For eksempel inneholder \rref{maalstk} en formel for 'målestokk'. Når de andre størrelsene er gitt, kan vi sette disse inn i formelen for å finne 'målestokk'. Slik kan vi si at vi da finner 'målestokk' \textsl{direkte}. Har man gjort oppgaver tilknyttet de tidligere kapitlene, har man allerede øvd rikelig på det å finne størrelser \textsl{direkte}.\vsk

I denne seksjonen skal vi se på det å finne størrelser \textsl{indirekte}. Med det mener vi at minst én av følgende gjelder:
\begin{itemize}
	\item Vi må løse en likning for å finne den ukjente størrelsen.
	\item Vi må ut fra en situasjonsbeskrivelse sette opp en formel som inneholder den ukjente størrelsen.
\end{itemize}


\info{Merk}{
I denne seksjonen er det bare gitt eksempler, og ingen regler. Det er fordi vi bruker regler vi har sett på i kapitlene om ligninger og funksjoner i \mb. Forskjellen er bare at vi her ser på størrelser med benevning. \vsk

På side \pageref{nondimension} så vi hvordan vi kan utføre utregninger uten å ta med benevningene til størrelsene. Det vil vi utelukkende gjøre i dette kapittelet, og derfor vil vi ikke her markere størrelser med \sym{*}.
}
\vsk
\eks[1]{
For en taxi er det følgende kostnader:
\begin{itemize}
	\item Du må betale 50\enh{kr} uansett hvor langt du blir kjørt.
	\item I tillegg betaler du 15\enh{kr} for hver kilometer du blir kjørt.
\end{itemize}
\abc{
\item Sett opp et uttrykk for hvor mye taxituren koster for hver kilometer du blir kjørt.
\item Hva koster en taxitur på 17\enh{km}?
}

\sv \vs
\abc{
\item Her er det to ukjente størrelser; 'kostnaden for taxituren' og 'antall kilometer kjørt'. Relasjonen mellom dem er denne:
	\[ \text{kostnaden for taxituren}=50+15\cdot\text{antall kilometer kjørt} \]
\item Vi har nå at
\[ \text{kostnaden for taxituren}=50+15\cdot17= 305 \]
Taxituren koster altså 305\enh{kr}.
}
}
\info{Tips}{
	Ved å la enkeltbokstaver representere størrelser, får man kortere uttrykk. La $ k $ stå for 'kostnad for taxituren' og $ x $ for 'antall kilometer kjørt'. Da blir uttrykket fra \textsl{Eksempel 1} over dette:
	\[ k=50+15x \]
	I tillegg kan man gjerne bruke skrivemåten for funksjoner:
	\[ k(x)=50+15x \] 
}
\eks[2]{
	Tenk at klassen ønsker å dra på en klassetur som til sammen koster 11\,000\,kr. For å dekke utgiftene har dere allerede skaffet 2\,000\,kr, resten skal skaffes gjennom loddsalg. For hvert lodd som selges, tjener dere 25\,kr.\os
	
	\abc{
		\item Lag en likning for hvor mange lodd klassen må selge for å få råd til klasseturen.
		\item Løs likningen.
	}
	
	\sv
	\abc{
		\item Vi starter med å tenke oss regnestykket i ord:
		\footnotesize
		\[ \text{penger allerede skaffet}+\text{penger per lodd}\cdot\text{antall lodd}=\text{prisen på turen} \]
		\normalsize
		Den eneste størrelsen vi ikke vet om er 'antall lodd'. Vi erstatter\footnote{Dette gjør vi bare fordi det da blir mindre for oss å skrive.} \textit{antall lodd} med $ x $, og setter verdiene til de andre størrelsene inn i likningen:
		\[ 2\,000+25\cdot x = 11\,000 \]
		\item \ \vs \vs
		\alg{
			25x &= 11\,000-2\,000\\
			25 x &= 9\,000\\
			\frac{\cancel{25} x}{\cancel{25}} &= \frac{9\,000}{25} \\
			x &= 360
		}
	}
}
\eks[3]{En vennegjeng ønsker å spleise på en bil som koster 50\,000 kr, men det er usikkert hvor mange personer som skal være med på å spleise.
	\abc{
	\item Kall 'antall personer som blir med på å spleise' for $ P $ og 'utgift per person' for $ U $,  og lag en formel for $ U $.
	\item Finn utgiften per person hvis 20 personer blir med.
	}
	\sv \vs
	\abc{
	\item Siden prisen på bilen skal deles på antall personer som er med i spleiselaget, må formelen bli
	\[ U = \frac{50\,000}{P} \]
	\item Vi erstatter $ P $ med 20, og får
	\[ U = \frac{50\,000}{20}= 2\,500 \]
	Utgiften per person er altså 2\,500\enh{kr}.
	}
}
\newpage
\eks[4]{
	Et sportsklubb planlegger en tur som krever bussreise. De får tilbud fra to busselskap:
	\begin{itemize}
		\item \textbf{Busselskap 1} \\
		Klassen betaler 10\,000\,kr uansett, og 10\enh{kr} per km.
		\item \textbf{Busselskap 2} \\
		Klassen betaler 4\,000\,kr uansett, og 30\enh{kr} per km.
	\end{itemize}
	For hvilken lengde kjørt tilbyr busselskapene samme pris?
	
	\sv
	
	Vi innfører følgende variabler:
	\begin{itemize}
		\item $ x=\text{antall kilometer kjørt} $ \\
		\item $ f(x)=\text{pris for Busselskap 1} $ \\
		\item $ g(x)=\text{pris for Busselskap 2} $
	\end{itemize}
	Da er \vs
	\alg{
		f(x)&=10x+10\,000\vn
		g(x)&=30x+4\,000
	}
	Busselskapene har samme pris når
	\alg{
		f(x)&=g(x) \\
		10x+10\,000&=30x+4\,000 \\
		6\,000&=20x \\
		x&=300
	}
	Busselskapene tilbyr altså samme pris hvis sportsklubben skal kjøre 300\enh{km}.
	%\fig{lig2}
}

\eks[5]{
	\textit{Ohms lov} sier at strømmen $ I $ gjennom en metallisk leder (med konstant temeperatur) er gitt ved formelen
	\[ I = \frac{U}{R} \]
	hvor $ U $ er spenningen og $ R $ er resistansen.
	
	\abc{
	\item Skriv om formelen til en formel for $ R $.	
}
Strøm måles i Ampere (A), spenning i Volt (V) og motstand i Ohm ($ \Omega $).
\abcs{2}{
\item Hvis strømmen er 2\,A og spenningen 12\,V, hva er da resistansen?
}

	
	\sv \vs
\abc{
	\item Vi gjør om formelen slik at $ R $ står alene på én side av likhetstegnet:\vs
\alg{
	I\cdot R&=\frac{U\cdot \cancel{R}}{\cancel{R}} \br
	I\cdot R &= U \br
	\frac{\cancel{I}\cdot R}{\cancel{I}} &= \frac{U}{I}\br 
	R &= \frac{U}{I}
}
\item Vi bruker formelen vi fant i a), og får at
\alg{
	R &= \frac{U}{I} \br
	&= \frac{12}{2} \\
	&= 6
}
Resistansen er altså $ 6\,\Omega $.
}
}
\newpage
\eks[6]{
	Gitt en temperatur $ T_C $ målt i antall grader Celsius ($ ^\circ$C). Temperaturen $ T_F $ målt i antall grader Fahrenheit ($ ^\circ$F) er da gitt ved formelen
	\[ T_F = \frac{9}{5}\cdot T_C+32 \]
\abc{
\item Skriv om formelen til en formel for $ T_C $.
\item Hvis en temperatur er målt til 59$ ^\circ F $, hva er da temperaturen målt i $ ^\circ C $?
}
	
	\sv \vs
	
	\abc{
	\item Vi isolerer $ T_C $ på én side av likhetstegnet:
	\alg{
		T_F &= \frac{9}{5}\cdot T_C+32 \\
		T_F-32 &= \frac{9}{5}\cdot T_C \\
		5(T_F-32) &= \cancel{5}\cdot\frac{9}{\cancel{5}}\cdot F_C \\
		5(T_F-32) &= 9T_C \\
		\frac{5(T_F-32)}{9} &= \frac{\cancel{9}T_C}{\cancel{9}} \\
		\frac{5(T_F-32)}{9} &= T_C
	}
	\item Vi bruker formelen fra a), og finner at
	\alg{
		T_C&= \frac{5(59-32)}{9} \br
		&= \frac{5(27)}{9} \br
		&= 5\cdot 3 \\
		&= 15
	}
59$ ^\circ $\enh{F} er altså det samme som 15$ ^\circ $\enh{C}. 
}
}
\section{Regresjon \label{Regresjon}}
Å forsøke å beskrive hvordan noe vil \textsl{utvikle} seg er en av de viktigste anvendelsene for funksjoner. Hvis vi har et datasett som beskriver tidligere hendelser, kan vi prøve å finne den funksjonen som passer best til datasettet. Dette kalles å utføre \outl{regresjon}. \vsk

Grafen under viser\footnote{Tall hentet fra \net{https://elbil.no/om-elbil/elbilstatistikk/}{elbil.no}} antall elbiler i Norge etter år 2010.
\fig{elbilsalg1}
Vi ønsker nå å finne en funksjon som
\begin{enumerate}[label=(\roman*)]
	\item så godt som mulig skjærer hvert punkt.
	\item har en graf som passer til situasjonen vi modellerer.
\end{enumerate}
Hvis vi utfører regresjon med en lineær funksjon i GeoGebra (se side \pageref{ggbreg}), får vi denne grafen:
\fig{elbilsalg2}
Utfører vi regresjon også med en andregradsfunksjon, får vi følgende resultat:
\fig{elbilsalg3}
I figuren over kan vi merke oss at
\begin{itemize}
	\item begge modellene (funksjonene) ''oppfører'' seg feilaktig i starten. Den lineære funksjonen starter med et negativt antall biler, mens den kvadratiske funksjonen starter med at antallet synker fra år 0 til år 1.
	\item Grafen til den kvadratiske passer punktene mye bedre enn grafen til den lineære funksjonen.
\end{itemize}
Hvis vi hadde antatt at den lineære funksjonen ga en god beskrivelse av antallet elbiler fremover i tid, kunne vi lest av fra grafen at antall elbiler i 2021 var ca. 350\,000. Hadde vi i stedet antatt det samme om den kvadratiske funksjonen, kunne vi lest av fra grafen at antall elbiler i 2021 var litt over 425\,000. Fasit er at antall elbiler i 2021 var 455\,271.
\fig{elbilsalg4}
\newpage
\section{Vekstfart, fart og akselerasjon}
\regdef[Gjennomsnittlig og momentan vekstfart]{
	Gitt en funksjon $ f(x) $. Da har vi at
	\begin{itemize}
		\item stigningstallet til linja som går gjennom punktene $ (a, f(a)) $ og $ (b, f(b)) $ kalles den \outl{gjennomsnittlige vekstfarten} til $ f $ på intervallet $ [a, b] $.
		
		\item $ f'(a) $ kalles den
		\outl{momentane vekstfarten} til $ f $ i $ a $. 
	\end{itemize}
}
\info{En praktisk tolkning av begrepene}{
	I \mb\ har vi sett at stigningstallet til linja som går gjennom punktene $ (a, f(a)) $ og $ (b, f(b)) $ er gitt ved uttrykket
	\[ \frac{f(b)-f(a)}{b-a} \]
	Hvis vi antar at dette forholdet er konstant for $ {x\in[a, b]} $, antar vi at $ f $ og $ x $ representerer proporsjonale størrelser på intervallet.\vsk
	
	Hvis vi ser tilbake til (den alternative) definisjonen av den deriverte i \tmen, innser vi at $ f'(a) $ er den gjennomsnittlige vekstfaktoren til $ f $ på intervallet $ [a, b] $ når $ {b\to a} $. Da går $ [a, b] $ mot å inneholde bare ett element, som er $ a $.
}
\newpage
\eks[]{
	Se for deg at vi slipper en ball fra 50 meter over bakken, og lar den falle fritt nedover. Når vi slipper ballen, starter vi også en stoppeklokke. Antall meter $ h $ ballen er over bakken etter $ t $ sekunder kan da tilnærmes ved funksjonen
	\[ h(t)=5\left(10-t^2\right)\quad,\quad t\in\left[0, \sqrt{10}\right] \]
	\abc{
		\item Finn den gjennomsnittlige vekstfarten til $ h $ på intervallet $ t\in[1, 3] $. Gi en praktisk tolkning av denne verdien.
		\item Finn den momentane vekstfarten til $ h $ i 3. Gi en praktisk tolkning av denne verdien.
	} \vs
	\sv
	\vs
	\abc{
		\item $ {h(1)=5\left(10-1\right)=45}  $ og $ {h(3)=5\left(10-3^2\right) =5} $. Det betyr at stigningstallet til linja mellom $ (1, h(1)) $ og $ (3, h(3)) $ er
		\[ \frac{5-45}{3-1}=-20 \]
		
		Altså er vekstfarten til $ h $ på intervallet $ [1, 3] $ lik $ -20 $. Siden $ h $ representerer antall meter, og $ t $ representerer antall sekunder, representerer vekstfarten en størrelse med enheten 'm/s'. Dette er en enhet for fart. Hvis ballen skulle falt 40 meter \textsl{nedover} i løpet av 2 sekunder med konstant fart, måtte denne farten vært $ 20\enh{m/s} $.
		\fig{ballb_am2}
		\item Siden $ {h'(t)=-10t }$, er $ {h'(3)=-30} $. Dette betyr at etter å ha falt i 3 sekunder, har ballen oppnådd farten $ 30\enh{m/s} $, i retning \textsl{nedover}.
	}
}
\reg[Farts- og akselerasjonsvektor \label{fartogaksvekt}]{
	Gitt en vektorfunksjon $ \vec{r}(t) $, hvor $ \vec{r} $ representerer en posisjon og $ t $ representerer tid. Da har vi at
	\begin{itemize}
		\item $ \vec{r}\,'(t) $ kalles \outl{fartsvektoren} og $ | \vec{r}\,'(t)|$ kalles \outl{banefarten}.
		\item $ \vec{r}\,''(t) $ kalles \outl{akselerasjonsvektoren} og	 $ | \vec{r}\,''(t)|$ kalles\\ \outl{baneakselerasjonen}.
	\end{itemize}
}
\fork{\ref{fartogaksvekt} Farts- og akselerasjonsvektor}{
	Hvis $ \vec{r}(t) $ representerer en posisjon (altså en relativ lengde fra et referansepunkt), og $ t $ en tid, vil $ \vec{r}\,'(t) $ innebære en lengde delt på en tid. Da vil $ \vec{r}\,'(t) $ representere en størrelse med en enhet for fart. $ \vec{r}\,''(t) $  vil innebære en fart delt på en tid, som da vil representere en akselerasjon.
}
\end{document}


